\section{Lattice libraries}\label{app:E.17}

Every lattice channel of \cref{chap:11} uses the same few objects for a given crystal: force constants, the frequencies and polarization vectors they give, the per-atom phonon density of states, the displacement tensors $B_\kappa$, and, in a polar crystal, the Born charges and $\varepsilon_\infty$. The multiphonon tables of \cref{sec:11.3} follow from them in minutes. For the candidate targets named in \cref{sec:17.5} most of the same numbers are not tabulated, and some of the tables need more than harmonic DFPT: all-electron densities at large momentum, path-integral momentum distributions, isotope-resolved Hessians and frequency-dependent Born charges. The DFPT method is that of \cref{app:E.4}, the anharmonic and path-integral methods those of \cref{app:E.8}.

\subsection{The shared harmonic base}

One harmonic lattice per material serves every lattice channel, and a mismatch between channels computed from different lattices is an error that no single channel reveals. The base consists of the force constants $\Phi_{\kappa\kappa'}(\Delta\vect{R})$, the eigenpairs $(\omega_{\vect{Q}\nu},\vect{e}_{\kappa\nu}(\vect{Q}))$, the per-atom density of states $D_\kappa(\omega)$, the tensors $B_\kappa$ and their off-site generalization $B_{\kappa\kappa'}(\Delta\vect{R})$ (\cref{eq:11.3-Bkk}), and for GaAs $Z^*_\kappa$ and $\varepsilon_\infty$. On the truncation label $\mathfrak{T}$ (\cref{sec:8.5}) it is at $n_B = 0$ and $n_A = 0$, static DFPT on the harmonic Born--Oppenheimer surface. The setting is the one of the representative calculations of \cref{app:E.1} and \cref{sec:5.2}, shared by every specification of \crefrange{chap:31}{chap:32}: PBE, SG15 norm-conserving pseudopotentials at \SI{60}{Ry}, and the measured room-temperature lattice constants of \cref{tab:E-functionals}, \SI{5.431}{\angstrom} (Si), \SI{5.658}{\angstrom} (Ge) and \SI{5.653}{\angstrom} (GaAs), with the GaAs frequencies also reported at the static-lattice constant \SI{5.638}{\angstrom}, at which \cref{sec:5.5} quotes them. The germanium and gallium pseudopotentials carry the 3d shell, which converges the total energy more slowly (\cref{app:E.1}), so for those two the frequencies are checked once at \SI{80}{Ry}. Silicon uses a $12^3$ $\vect{k}$ grid and an $8^3$ wavevector grid (29 irreducible points), which the phonon-assisted window needs (\cref{app:E.14}); germanium and GaAs use $8^3$ and $6^3$. DFPT runs to a self-consistency threshold of $10^{-16}$, with the electric-field perturbation for GaAs, the potential changes kept on disk for the vertices, and Fourier interpolation with the acoustic sum rule imposed on the force constants. GaAs needs the nonanalytic term of \cref{eq:E.4-na}, without which its longitudinal-optical branch is wrong at small $\vect{Q}$.

One low-symmetry wavevector of silicon at \SI{40}{Ry} with an $8^3$ $\vect{k}$ grid takes about eight minutes on four cores; scaled by $(60/40)^{3/2}$ for the cutoff and $(12/8)^3$ for the grid, the 29 wavevectors take about \SI{24}{h} on four cores, roughly 100 core-hours, in less than \SI{2}{GB}. Germanium takes 13 to 16 minutes per wavevector on four cores at its smaller grids, about 15 core-hours, and GaAs with the electric field about 20.

The frequencies change by less than \SI{0.1}{meV} between the two coarsest wavevector grids, and they are compared with neutron data at $\mathrm{\Gamma}$, X and L (silicon~\cite{nilsson1972study}, germanium~\cite{nilsson1971phonon}) with a tolerance per branch, since the error of PBE at the measured lattice constant differs from branch to branch. The optical and longitudinal-acoustic branches pass at \SI{3}{\percent}. The transverse-acoustic branches of silicon are \SI{8}{\percent} soft at X and L in this setting (\cref{sec:5.2}), and pass at \SI{10}{\percent}; those of germanium are 11.6 to \SI{13.5}{\percent} soft on the $4^3$ grid and pass at \SI{15}{\percent} until the test of \cref{app:E.14} has separated the grid error from the functional. Where a transverse-acoustic frequency sets an edge, the neutron-scaled branch is the reference and the unscaled one is reported beside it. The phase convention matters for every coherent quantity downstream: the eigenvector carries the cell phase $\ee^{\ii\vect{Q}\cdot\vect{R}}$, and the intracell phase $\ee^{\ii\vect{q}\cdot\vect{\tau}_\kappa}$ is taken at the unreduced $\vect{q}$ (\cref{app:A.3}).

The deformation-potential tensors $D^n$ of \cref{eq:6.2-D} belong with the base, since they set the soft acoustic limit of the phonon-assisted amplitude (\cref{sec:13.6,app:D.37}). They are computed by frozen strain: each of the six strain components is applied at $\pm\SI{0.5}{\percent}$ to the two-atom cell with the internal coordinate relaxed, and the band energies at the conduction minimum and at $\mathrm{\Gamma}$ are referenced to the average electrostatic potential. Only differences between the conduction and valence tensors enter at the band extrema, so the absolute alignment from a slab calculation is needed only for single-band checks against published deformation potentials~\cite{li2021deformation}. At the extrema the group velocity vanishes, and the comoving intraband vertex of \cref{eq:D.37-Xi} reduces to the symmetric tensor that frozen strain computes; away from them it carries the Galilean term $\vect{v}^{\mathrm{g}}_n\otimes\vect{p}_{nn}$, which frozen strain does not produce and which is added from the band velocities and momentum matrix elements. Valley degeneracy is lifted under shear strain, and the bands must be followed by character, not by index. The run takes minutes, and its size decides which terms dominate the soft limit: with deformation potentials of order \SI{9}{eV} and a longitudinal sound speed of \SI{8.43}{km.s^{-1}}, the ratio $D/c_\nu$ is 300 to \SI{350}{keV} (\cref{sec:13.6}), far above the momentum transfers of a few keV, so the deformation terms dominate the soft limit by two orders of magnitude.

\subsection{Multiphonon tables and coherence}

At momentum transfers far above the zone boundary a nuclear recoil excites many phonons, and in the incoherent approximation of \cref{sec:11.3} the response of each atom follows from its own density of states. On a cubic site the directional density of states is a third of $D_\kappa(\omega)$, so the per-atom density of states is the complete input. The library holds $S(\vect{q},\hat{\vect{q}},\omega)$ per cell, in \si{eV^{-1}}; the per-volume $S$ of \cref{sec:4.2}, in \si{eV^2} in natural units, is this divided by the cell volume $\Omega$. The library also holds the probability $P_{\mathrm{ph}}(n|\vect{q},\omega)$ that $n$ phonons were created, the branch content per atom, the coherent fraction and the coherence correction below, on a grid of 400 logarithmic points in $q$ from 0.1 to \SI{300}{keV}, $\omega$ from 0 to \SI{1}{eV} in steps of \SI{0.1}{meV}, and temperatures of \SI{10}{mK}, \SI{1}{K} and \SI{77}{K}. Deep in the impulse regime, for $x_R > 16$, the table switches to the impulse Gaussian with the third- and fourth-cumulant corrections whose inputs come from the path-integral calculation below. The level is $\mathcal{K}$ = multiphonon, incoherent, with the coherent one-phonon term added and $n_A = 0$.

At zero temperature the multiphonon expansion is a forward recursion in $n$; at finite temperature the time-domain form $\ee^{-2W_\kappa}[\ee^{\mathcal{C}_{\kappa\kappa}(t)} - 1]$ of \cref{sec:5.6} is Fourier transformed with a Gaussian regulator in time. The densities of states come from $48^3$ grids with tetrahedra. The code is about 300 lines on top of the base and runs in minutes. Its validation is an Einstein site, whose $n$-phonon weights are known in closed form at any temperature and are reproduced to $10^{-5}$. The convergence tests are the $f$-sum rule of the lattice, $\int\omega S\,\dd\omega$ fixed by the recoil energy, to $10^{-4}$ at every $q$, and a time window $t_{\mathrm{max}} \ge 20\pi/\delta\omega$ for the bin width $\delta\omega$. The exponent $2W_\kappa$ uses $B_\kappa$ with the thermal factor $2n_{\vect{Q}\nu} + 1$; the elastic Bragg peak is removed analytically; and GaAs needs a separate $W_\kappa$ for each species. When the Migdal channel is computed separately, the lattice response is multiplied by $1 - P_{\mathrm{Mig}}$ and uses the dressed mass.

Below the zone boundary the incoherent approximation fails pointwise for two-phonon emission, because the two phonons from neighboring atoms interfere. The coherent two-phonon response is a Brillouin-zone sum over one phonon wavevector, the second fixed by momentum conservation, and its $\omega$ integral has an exact real-space form. With the static correlation tensor $B_{\kappa\kappa'}(\Delta\vect{R})$ of \cref{eq:11.3-Bkk}, the $n$-phonon weight is the lattice sum \cref{eq:11.3-weight}, per cell, and its ratio $\delta^{\mathrm{coh}}_n(q)$ to the incoherent weight is the exact coherence correction for bins that cover the whole $n$-phonon band. At zero temperature in three dimensions $B_{\kappa\kappa'}$ falls as $\Delta R^{-2}$, so the sum converges absolutely for $n \ge 2$. The lattice sum \cref{eq:11.3-weight} is not used at finite temperature, where the thermal factor spoils the convergence, nor at $n = 1$; at \SI{1}{K} and \SI{77}{K} the coherence correction comes from the pointwise zone sum below, which converges at any temperature. The pointwise response $S^{(2)}_{\mathrm{coh}}(\vect{q},\omega)$ is computed for silicon, then germanium and GaAs, on 40 logarithmic values of $|\vect{q}|$ from 0.2 to \SI{10}{keV}, the 20 or so irreducible directions of a 590-point Lebedev grid, and bins of \SI{0.5}{meV}, with $B_{\kappa\kappa'}$ to \SI{40}{\angstrom} by FFT from $48^3$ eigenpairs and the zone sum on $48^3$ with a Gaussian of \SI{0.25}{meV}. The tests are a change below \SI{2}{\percent} per bin from $32^3$ through $48^3$ to $64^3$ and below $10^{-3}$ in $\delta^{\mathrm{coh}}_n$ under the real-space radius. The acoustic point $\vect{Q} = 0$ has an integrable singularity that a mesh samples badly: it is excluded and a small sphere around it is added analytically. The whole run costs a few core-hours after the base.

The same tables, projected on cubic harmonics, give the anisotropy of the lattice response for \cref{sec:11.7}. In the harmonic incoherent part a cubic site is isotropic at every phonon order, so the components with $\ell \ge 4$ vanish there to numerical noise, below $10^{-6}$ of the isotropic part; a larger value is a symmetry-breaking artifact of the mesh. The anisotropy comes only from the coherent terms, chiefly the one-phonon term. At long wavelength its $\ell = 4$ component in silicon is $S_4/S_0 = 0.028$ with the measured elastic constants and 0.025 with the PBE force constants at $q \le 0.3\,q_{\mathrm{BZ}}$, and $|S_6/S_0|$ is below $10^{-3}$ (\cref{sec:11.7}); the components are tabulated as $S_\ell/S_0$ on the cubic harmonics $Y^{\mathrm{c}}_\ell$. The projection uses the 590-point Lebedev grid, exact to degree 41, with $\ell \le 12$ for cubic crystals and $\ell \le 8$ for sapphire and SiC, and $\omega$ bins wider than the change of the dispersion between neighboring directions.

\subsection{The charge vertex at large momentum}

Two corrections to the lattice vertex need all-electron densities. The first is the valence charge deformation of \cref{sec:11.3}: whether the electron cloud follows a displaced nucleus at $q$ far above the zone boundary. Split into a rigid neutral atom and a deformation, the charge vertex of atom $\kappa$ is $q_\alpha[Z_\kappa - f^{\mathrm{at}}_\kappa(q)] - \Delta_{\kappa\alpha}(\vect{q})$, with $f^{\mathrm{at}}_\kappa$ the all-electron atomic form factor, and the deformations obey an exact sum rule: summed over the atoms of the cell at $\vect{Q} = 0$, the density changes give $G_\alpha$ times the Fourier component of the electron density at $\vect{G}$, the response to a uniform translation. The deformation $\Delta_{\kappa\alpha}(\vect{Q}+\vect{G})$ is wanted for Si, Ge and GaAs at $|\vect{Q}+\vect{G}| \le \SI{20}{keV}$, about 725 reciprocal-lattice vectors per $\vect{Q}$ in silicon and 821 in germanium and GaAs. Up to about \SI{10}{keV} it is the DFPT density change of the base, Fourier transformed per $\vect{Q}$, minus $q f^{\mathrm{at}}_\kappa$. At all $q$ it is the symmetric difference of all-electron densities with one atom displaced by $\pm\SI{0.01}{\angstrom}$ in a 216-atom cell, from PAW with all-electron reconstruction or from a full-potential code. The PAW density on the FFT grid aliases its core-region components above about \SI{10}{keV}, so those components are taken as the smooth-part FFT plus one-center radial transforms per atom. The density cutoff must reach 1.5 times the largest $q$, at least \SI{65}{Ry}, or what the PAW dataset needs if that is more. The result must satisfy the sum rule to $10^{-3}$ and be linear in $\pm u$. A frozen-core calculation gives no core polarization; that part, the polarization of the core by the residual field at the moving nucleus, comes from a Sternheimer (Dalgarno--Lewis) calculation on the all-electron core orbitals, checked by the sum rule that its dipole moment equals the core polarizability. The DFPT part costs nothing extra, the supercells about $10^2$ core-hours, the cores minutes.

The second correction is the finite-$q$ factor of the electron density at the nucleus (\cref{sec:12.5}), which enters the cross term between nuclear and electronic recoil. Its spherical atomic part is known from all-electron atoms; the crystal anisotropy and the relativistic enhancement are not. The crystal valence adds an anisotropy of tetrahedral symmetry, estimated at $1.3\times10^{-3}$ of the factor in silicon; it is computed from the valence all-electron density (PAW minus the atomic core, or a full-potential code) to $|\vect{G}| \le 40/a_0$ on a Lebedev grid, converged to $10^{-5}$ of the density at the nucleus under the cutoff and the augmentation radius. The relativistic density averaged over the finite nucleus raises the contact density by the factor $\mathcal{R}_{\mathrm{rel}}$ of \cref{eq:12.5-Rrel}, \SI{8}{\percent} in silicon and 40 to \SI{45}{\percent} in Ge, Ga and As. It comes from a finite-nucleus Dirac calculation of the atom with a Fermi or uniform nuclear charge of the measured radius, with the innermost radial point below a hundredth of the nuclear radius; the choice of nuclear model changes it by less than \SI{1}{\percent}. A relativistic point-nucleus density diverges at the origin, and its value at the first grid point is an artifact of the grid. Norm-conserving pseudo-densities are unusable here. The crystal part costs about 10 core-hours, each atom minutes.

\subsection{Nuclear momentum distributions}

The width and skewness of the impulse peak of \cref{sec:11.4} are set by the quantum motion of the nuclei, which harmonic theory gives only to leading order. Path-integral molecular dynamics with i-PI~\cite{kapil2019ipi}, driving a machine-learned potential (\cref{app:E.11}) trained on PBE or SCAN (strongly constrained and appropriately normed) forces in 216-atom cells, gives the momentum distribution of each species to all orders in the anharmonicity: the centroid-virial estimator yields the kinetic energy along $\hat{\vect{q}}$, the open-path estimator the momentum distribution and the fourth moment, and the analytic Hessian averaged over beads the mean curvature of the Born--Oppenheimer surface, which is the third cumulant of the impulse peak. At \SI{10}{K} silicon would need about 250 beads, from $\omega_{\mathrm{max}}/(k_BTP) \lesssim 0.3$, so we run at 50 to \SI{100}{K} with a colored-noise thermostat, which brings the bead number to 16 to 32, and correct to zero temperature with the harmonic difference from the base. The kinetic energy anchors the anharmonic template, and the anharmonic change of the kinetic energy of silicon is at most $\pm\SI{0.2}{\percent}$ (\cref{sec:11.5}); its error, statistical plus that of the harmonic-difference correction, must therefore be at most \SI{0.2}{\percent}. The potential is validated against DFT on displaced configurations out to three standard deviations of the sampled distribution, since extrapolation beyond the training set biases the fourth cumulant. The open-path estimator of the fourth moment is noisy and needs at least $10^6$ steps, and the whole calculation costs about $10^3$ core-hours per material.

\subsection{Isotope-resolved Hessians}

A vector coupled to baryon minus lepton number couples to the isotope disorder of the lattice (\cref{sec:11.6}), and its absorption rate $\varGamma_{\mathrm{iso}}(\omega)$ is computed exactly by diagonalizing the Hessian of a disordered supercell. For natural silicon the supercell has 4096 atoms ($8\times8\times8$ conventional cells, 12{,}288 degrees of freedom, about ten minutes per dense diagonalization), with the force constants of the base and ten random isotope configurations; the vertex is the mass-subtracted one, with the mean taken mass-weighted, otherwise a spurious acoustic term appears. The check is the exact single-site (Lifshitz) Green function, which gives the leading correlation correction between vertex and eigenvector; with the PBE density of states it is $+\SI{8.0}{\percent}$ at \SI{50}{meV} (\cref{sec:11.6}). The rate must change by less than \SI{2}{\percent} from 10 to 20 configurations and from 4096 to 13{,}824 atoms, the larger cell by the kernel polynomial method; the whole calculation costs about 10 core-hours.

The competing electromagnetic absorption at the same frequency is two-phonon absorption, and the bound on the ratio in \cref{sec:11.6} uses room-temperature data where the millikelvin spectrum is needed. The zero-temperature two-phonon $\varepsilon_2(\omega)$ of silicon follows from the second-order dipole moments $\partial Z^*/\partial u$, by the first-principles route of Deinzer and Strauch~\cite{deinzer2004two}. The derivatives come from finite differences of DFPT Born charges in displaced 64- and 216-atom cells, with displacements of $\pm0.005$, $\pm0.01$ and $\pm\SI{0.02}{\angstrom}$ and linearity to \SI{2}{\percent} (in silicon $Z^* = 0$ at equilibrium, so the first order is linear in $u$), at \SI{80}{Ry}, with $4^3$ and $2^3$ $\vect{k}$ grids checked against $6^3$ and $3^3$. The two-phonon sum uses only the sum processes, which survive at zero temperature, on $32^3$ and $48^3$ grids with Gaussians of 0.25 to \SI{1}{meV} extrapolated to zero. The test comes first: the same calculation with room-temperature occupations, sum and difference processes, and quasi-harmonic frequencies reproduces the measured room-temperature absorption to \SI{20}{\percent} between 40 and \SI{60}{meV}, and only then is the zero-temperature curve used. The two cells cost about $10^2$ and $2\times10^3$ core-hours.

\subsection{Dynamical Born charges}

A recoiling ion in a polar crystal radiates through its Born charge at the photon frequency, not at zero frequency (\cref{sec:12.6}). The calculation needs $Z^*_\kappa(\omega)$ at zero macroscopic field for $\omega$ from \SI{0.05}{eV} to the gap, with the refractive index. At small $\omega$, $Z^*_\kappa(\omega) = Z^*_\kappa(0) - \omega^2C_\kappa + O(\omega^4)$, and a Cauchy--Schwarz bound on the same spectral sums gives for GaAs $|Z^*(\omega) - Z^*(0)| \le 0.9\,(\omega/\SI{1}{eV})^2$ against $Z^*(0) \approx 2.1$: at photon energies near \SI{1}{eV} the dispersion is not negligible a priori and must be computed. Frequency-dependent DFPT computes it as a mixed response to an electric field and an atomic displacement, solved by shifted Sternheimer equations (\cref{eq:E.4-sthw}), at the RPA and then with the adiabatic local-density kernel. The PBE gap would place the resonance of $Z^*(\omega)$ far below the true edge, so the bands carry a scissor to the measured gap, and the response is taken at zero macroscopic field, since with the field the photon line is counted twice. Imaginary frequencies give $C_\kappa$ by fits in $\omega^2$; real frequencies near the gap use a damping of 10 to \SI{20}{meV}, extrapolated to zero. $Z^*(0)$ and $\varepsilon_\infty$ must agree with static DFPT to \SI{1}{\percent}, and $\sum_\kappa Z^*_\kappa(\omega)$ must tend to 0 as $\omega\to0$ and to $\sum_\kappa Z_\kappa$ as $\omega\to\infty$. Each frequency batch costs about 50 core-hours for GaAs.

For a scalar coupled to electrons, the single-phonon rate of GaAs is a near-cancellation (\cref{sec:11.8}). The vertex enters through $d = v_{\mathrm{Ga}}/M_{\mathrm{Ga}} - v_{\mathrm{As}}/M_{\mathrm{As}}$ with $v_\kappa = Z_\kappa - Z^*_\kappa/\varepsilon_\infty$, which is $-1.21\times10^{-3}$ per u for $Z^* = 2.11$ and $\varepsilon_\infty = 10.88$, because the nuclear charges per unit mass of gallium and arsenic nearly coincide (\cref{sec:11.8}). A \SI{10}{\percent} change in $Z^*/\varepsilon_\infty$ changes $d$ by \SI{44}{\percent}, so a rate good to \SI{20}{\percent} needs $Z^*/\varepsilon_\infty$ to \SI{2.3}{\percent}. PBE-DFPT cannot deliver this, since its $\varepsilon_\infty$ for GaAs is tens of percent too large. The combination is taken instead from measured quantities through the Lyddane--Sachs--Teller relation~\cite{lyddane1941polar},
\begin{equation}
  \frac{Z^*}{\varepsilon_\infty} = \Big[\frac{\Omega\,\mu_{\mathrm{red}}(\omega_{\mathrm{L}}^2 - \omega_{\mathrm{T}}^2)}{e^2\varepsilon_\infty}\Big]^{1/2} ,
  \label{eq:E.17-lst}
\end{equation}
in Heaviside--Lorentz units ($e^2 = 4\pi\alpha$), with $\mu_{\mathrm{red}}$ the reduced mass of the pair and $\omega_{\mathrm{L}}$, $\omega_{\mathrm{T}}$ the zone-center longitudinal and transverse optical frequencies, all at low temperature, since $\varepsilon_\infty$ changes by 2 to \SI{3}{\percent} between room temperature and \SI{4}{K}. With $\varepsilon_\infty$ to \SI{1}{\percent} and the Raman frequencies to \SI{0.1}{\percent} the requirement is met; DFPT supplies the sign of $Z^*_{\mathrm{Ga}}$ and a consistency check. The phonon widths come from three-phonon calculations, to \SI{10}{\percent}. Rigid-ion charges change $d$ by a factor of three and must not be used, and the gallium isotopes add a broadband channel of the isotope type above.

\subsection{Data for the candidate targets}

The comparison of \cref{sec:17.5} and the target table of \cref{sec:21.7} need, for each candidate, $\varepsilon_\infty$ and $\varepsilon_0$, the Born-charge tensors, the zone-center transverse and longitudinal optical frequencies, the sound speeds along the principal directions, $\omega_{\mathrm{max}}$, the phonon density of states, the elastic tensor, phonon lifetimes and the thermal contraction to millikelvin. Where a measured value exists, it is used; DFPT fills the rest. Each candidate gets a relaxation, a zone-center DFPT run with the electric field ($\varepsilon_\infty$, $Z^*$, the optical modes and, through the mode sum, $\varepsilon_0$), and $4^3$ and $6^3$ wavevector grids for the slopes, sound speeds, $\omega_{\mathrm{max}}$ and density of states. \Cref{tab:E17-batch} lists the cutoffs, set where the total energy converges to \SI{1}{meV} per atom and the pressure to \SI{1}{kbar}.

\begin{table}[t]
  \centering
  \caption{Cutoffs and costs of the zone-center and phonon-grid DFPT runs for the candidates beyond the shared base, with SG15 norm-conserving pseudopotentials. Costs are in core-hours; on one accelerator each small-cell run takes minutes to an hour. The last four rows are batches.}
  \label{tab:E17-batch}
  \small
  \begin{tabular}{llll}
    \toprule
    material & cutoff (Ry) & cell & cost (core-hours) \\
    \midrule
    diamond, 3C-SiC, LiF, NaF & 90--100 & 2 atoms & 2--20 \\
    4H-SiC & 90 & 8 atoms & 40--80 \\
    Al$_2$O$_3$ & 100 & 10 atoms & 120--240 \\
    CaWO$_4$, Li$_2$MoO$_4$ & W, Mo semicore & primitive cells & \num{1.2e3}--\num{5.4e3} \\
    magnetic, DFT $+\,U$ (YIG dominant) & per material & 4--80 atoms & \num{1.1e4}--$10^5$ \\
    23 further candidates & per material & small cells & $10^3$--\num{5e3} \\
    thermal contraction, five volumes & as above & all candidates & \num{5e3}--\num{2.5e4} \\
    \bottomrule
  \end{tabular}
\end{table}

PBE overestimates $\varepsilon_\infty$ systematically, by \SI{11}{\percent} in silicon (\cref{app:E.4}), so measured dielectric constants are used for ranking and the computed ones are diagnostics. The Lyddane--Sachs--Teller relation, $\varepsilon_0/\varepsilon_\infty = \prod_j\omega_{\mathrm{L},j}^2/\omega_{\mathrm{T},j}^2$, holds per principal axis and checks each run; a non-cubic crystal has a pair of dielectric constants per axis. In soft-mode crystals (SrTiO$_3$, the lead chalcogenides) the DFPT $\varepsilon_0$ is unusable, and the measured low-temperature value replaces it; an infrared mode below \SI{20}{cm^{-1}} flags such a case. Heavy elements (Pb, Bi, Te, W) need spin--orbit coupling at least as a check. The scintillating crystals CaWO$_4$ and Li$_2$MoO$_4$ are treated like the oxides, with the tungsten or molybdenum semicore in the valence and a measured elastic tensor where one exists. For the magnetic candidates (YIG, NiO, MnO, Cr$_2$O$_3$, NiPS$_3$) DFPT runs with a Hubbard $U$ from linear response~\cite{timrov2018hubbard} on the magnetic ground state, with antiferromagnetic cell doubling for NiO and MnO; the dependence of $\varepsilon$ and $Z^*$ on $U$ is reported as the uncertainty, and measured moments and infrared modes are the checks.

Elastic tensors come from strain--stress calculations or DFPT. Each is positive definite and in the Voigt order; at least one public database lists the shear components in a different order, a trap when tables are merged. For a cubic crystal the longitudinal sound speed in any direction is at least $\sqrt{B/\rho_{\mathrm{T}}}$ (\cref{sec:11.1}), which tests both computed and quoted values. Where a ranking depends on a computed tensor, as for TeO$_2$ against the measured CaWO$_4$ in \cref{sec:21.7}, the measured tensor is the requirement and the offset of computed tensors, \SI{4.9}{\percent} on average for silicon, germanium and diamond (\cref{sec:17.6-correlation}), is the uncertainty until it exists. CsI, whose sound speeds set its nuclear-recoil edge, enters with its measured elastic constants at \SI{4.2}{K}, $C_{11} = 27.37$, $C_{12} = 7.93$ and $C_{44} = \SI{8.25}{GPa}$; the thermal contraction below supplies the density of its cold cell.

Phonon lifetimes come from three-phonon perturbation theory, on-shell, which is $n_A = 2$ on the truncation label: third-order force constants from phono3py~\cite{togo2015distributions} on $3\times3\times3$ supercells, or from a machine-learned potential, converged in supercell size and interaction cutoff. The anharmonic decay constant of silicon~\cite{tamura1985spontaneous} is the check; for germanium the published values spread by a factor of ten. The cost is $10^3$ to $10^4$ core-hours per material, and the lifetimes enter \cref{sec:11.8}. The thermal contraction integrated from \SI{300}{K} to millikelvin, which stores elastic energy between a crystal and its holder, comes from the quasi-harmonic approximation (\cref{eq:E.8-qha}) with DFPT at about five volumes, with the PBE and PBEsol spread as the uncertainty and measured expansion coefficients (Si, Ge, sapphire) as checks. The quasi-harmonic approximation fails near soft modes, and silicon expands on cooling below about \SI{120}{K}, so the sign of the integrand must be checked.
